\subsection{Sprint 2}

Após uma boa entrega na Sprint 1, mantivemos as squads e eu decidi continuar no \textit{frontend}. Dessa vez, fiquei responsável por criar um componente de anexo de arquivos no aplicativo. Comecei a sprint estudando mais sobre React Native e procurando alternativas para permitir que o usuário selecionasse documentos dentro do fluxo da aplicação.

Inicialmente tentei utilizar a biblioteca \textit{Document Picker}, mas encontrei problemas de compatibilidade com o Expo. Essa dificuldade me travou por alguns momentos, porque eu ainda não tinha experiência suficiente para diferenciar um erro de implementação de uma limitação da biblioteca escolhida. Depois de pesquisar e conversar com colegas AGES III, entendi que o caminho mais adequado era usar o \textit{expo-document-picker}, que se integrava melhor ao projeto e as dependências já utilizadas. Com essa troca, consegui finalizar o componente de anexo de arquivos. Por conta da dificuldade inicial, fiquei apenas com uma entrega parcial, mas consegui implementar a funcionalidade principal e deixá-la pronta para integração com o restante do aplicativo. 

Com essa troca, consegui finalizar o componente de anexo de arquivos. A principal lição da sprint foi entender a importância de comunicar problemas cedo. Quando compartilhei a dificuldade, consegui receber direcionamento e evitar perder ainda mais tempo em uma solução que não se encaixava bem no ambiente do projeto. Mesmo 
